\section{Introduction}

The scratch-wound assay is a widely used \emph{in vitro} method for studying collective cell migration. By creating a cell-free gap within a confluent cell monolayer, time-lapse microscopy captures the progressive closure of the wound as cells at the wound margin proliferate and migrate inward \citep{Liang2007,Jonkman2014,Grada2017,Yarrow2004,Ashby2012,Pinto2023,Stamm2016}. Cell migration governs tissue repair, morphogenesis, and tumor invasion, making reliable quantification of wound closure dynamics a direct readout of migration capacity across pharmacological screening, developmental biology, and oncology.

Despite the assay's conceptual simplicity, quantification remains a significant bottleneck. Manual measurements are labor-intensive, introduce inter-observer variability, and do not scale efficiently to large experiments. To address these limitations, several computational tools have been developed, including TScratch \citep{Gebaeck2009}, ImageJ plugins \citep{Suarez2020}, the Robust Quantitative Scratch Assay \citep{Vargas2016}, AIM \citep{Nunes2017}, CellProfiler pipelines \citep{McQuin2018}, and more recently level-set methods with spatial sector analysis \citep{Martinotti2025}.

Scratch-wound experiments are inherently temporal, generating image sequences that capture the continuous progression of wound closure. While existing approaches have substantially improved throughput and reproducibility, segmentation is often performed on individual frames, with temporal information incorporated primarily during downstream analysis. Integrating temporal information directly into the segmentation stage enforces physical constraints on boundary evolution, eliminating the non-biological area fluctuations that accumulate when each frame is analyzed in isolation. Temporal constraints have previously been explored in related domains such as cell tracking \citep{Bragantini2025,Lee2023,Li2008} and video segmentation \citep{Guan2021}.

In this study, we propose an automated end-to-end pipeline for scratch-wound analysis that combines texture-based wound detection with temporal integration across image sequences. The pipeline uses local variance information to distinguish the wound channel from the surrounding cellular monolayer in phase-contrast images and incorporates temporal continuity to generate consistent wound-closure trajectories \citep{Yin2012,Treloar2013}. The pipeline additionally provides automated quality control, parallel processing of image sequences, and extraction of quantitative wound-closure metrics including wound area, closure speed, edge displacement, and closure fraction. The pipeline is fully automated and requires neither training data nor user-defined threshold tuning. To our knowledge, this is the first application of a hard monotonic closure constraint to scratch-wound segmentation.

We evaluated the pipeline using a standard scratch-wound migration assay in immortalized skeletal muscle cells \citep{Goetsch2011,Goetsch2013} imaged under three culture conditions: an ALK5i, a DMSO vehicle control, and a media-only control. These conditions provide a controlled test bed for assessing performance across distinct wound-closure trajectories.


%
%A variance-based single-frame detector that exploits the fundamental texture difference between cell monolayer (high local variance due to cell edges) and the wound channel (smooth, low variance). This approach is grounded in the physics of phase-contrast image formation \citep{Yin2012} and follows the precedent of training-free texture-based segmentation in label-free microscopy \citep{Ambrossio2022,Vicar2019}. No training data or threshold tuning by the user is required.

%
%Single-frame independence. Most tools process each image in isolation, ignoring the temporal structure of the experiment
%and producing noisy area-versus-time curves with non-biological fluctuations ---
%apparent wound re-opening that contradicts the physical reality of the assay
%\citep{Treloar2013}.
%This is particularly problematic when computing derived metrics such as closure speed, where non-monotonic noise propagates into erroneous rate estimates.
%
%Machine-learning approaches, including U-Net-based architectures
%\citep{Ronneberger2015,Dogru2024,Naser2025}
%and foundation models such as the Segment Anything Model \citep{Lehmann2024},
%require annotated ground-truth masks.
%These annotations are themselves subject to observer bias and are expensive to produce at scale.
%Deep-learning pipelines further introduce domain-shift risks when applied to imaging conditions
%different from their training distribution \citep{Oldenburg2021}.
%
%
%To our knowledge, no existing tool enforces the physical axiom that a scratch wound is expected to close over time. While post-hoc smoothing can reduce noise, it does not encode the underlying biological axiom and may introduce artifacts of its own.
%
%
%To evaluate the pipeline we use a standard scratch-wound migration assay in immortalized skeletal muscle cells \citep{Goetsch2011,Goetsch2013}, imaged under three culture conditions (an ALK5/TGF-$\beta$ receptor~I inhibitor, a DMSO vehicle, and a media control). These conditions serve solely as a controlled test bed: they provide trajectories that may or may not differ in closure dynamics, allowing us to ask whether the method recovers condition-dependent differences. This work makes no claim about the underlying pharmacology, and we draw no biological conclusions from the treatment comparison.
%
%
%
%We introduce \textbf{Quantification}, an automated two-stage pipeline that addresses the limitations of existing wound quantification tools.  A variance-based single-frame detector that exploits the fundamental texture difference between cell monolayer (high local variance due to cell edges) and the wound channel (smooth, low variance). This approach is grounded in the physics of phase-contrast image formation \citep{Yin2012} and follows the precedent of training-free texture-based segmentation in label-free microscopy \citep{Ambrossio2022,Vicar2019}. No training data or threshold tuning by the user is required.
%
%A stateful temporal integrator that encodes the monotonic closure axiom directly into the segmentation. Once a pixel column has been declared healed, it cannot revert. This eliminates non-biological noise in the area-versus-time curve and produces smooth, strictly non-increasing wound trajectories. While temporal constraints have been explored in cell tracking \citep{Bragantini2025,Lee2023,Li2008} and video segmentation \citep{Guan2021}, to our knowledge this is the first application of a hard monotonic constraint to wound-healing segmentation.
%
%The pipeline is complemented by automated multi-level quality control, parallel trajectory processing, and downstream metric extraction (closure speed, edge displacement, closing fraction). We validate Quantification on 124 wound-bearing phase-contrast images from across the three culture condition.
